%% ============================================================
%%  Chapter 2: Background and Related Work
%% ============================================================
\chapter{Background and Related Work}
\label{chap:background}

This chapter reviews the literature relevant to the two core problems addressed
in this thesis: multi-robot task allocation (Section~\ref{sec:bg_mrta}) and
multi-agent path finding (Section~\ref{sec:bg_mapf}).
Section~\ref{sec:bg_integrated} covers work that attempts to integrate the two,
Section~\ref{sec:bg_communication} reviews coordination under communication
constraints, and Section~\ref{sec:bg_warehouse} situates the work in the
context of warehouse automation.
Section~\ref{sec:bg_positioning} identifies the specific gap that this thesis
fills.


%% ── 2.1 ─────────────────────────────────────────────────────────
\section{Multi-Robot Task Allocation}
\label{sec:bg_mrta}

Multi-robot task allocation (MRTA) is the problem of deciding which robot
performs which task, and in what order, so as to optimize a system-level
objective such as makespan, total distance, or throughput.
The canonical taxonomy of \citet{gerkey2004formal} classifies MRTA problems
along three axes: \emph{single-task} (ST) versus \emph{multi-task} (MT) robots,
\emph{single-robot} (SR) versus \emph{multi-robot} (MR) tasks, and
\emph{instantaneous assignment} (IA) versus \emph{time-extended assignment}
(TA).
The warehouse problem studied in this thesis is an ST-SR-TA instance: each
robot handles one task at a time, each task requires exactly one robot, but
assignments evolve over time as new tasks arrive.

In its most general form, MRTA is NP-hard, even for single-task
robots~\cite{gerkey2004formal}.
Practical systems therefore rely either on polynomial-time approximation
algorithms with bounded sub-optimality, or on heuristic methods that sacrifice
formal guarantees for scalability.
The choice between centralized and decentralized architectures
      (Sections~\ref{subsec:bg_centralized} and~\ref{subsec:bg_decentralized})
is the primary design tension.


%% ── 2.1.1 ──────────────────────────────────────────────────────
\subsection{Centralized Approaches}
      \label{subsec:bg_centralized}

Centralized MRTA methods aggregate all information at a single planner that
computes assignments for all the agents.
The classical assignment problem (minimize total cost for a one-to-one
agent-task matching) can be solved optimally in $O(n^3)$ time by the Hungarian
algorithm~\cite{kuhn1955hungarian}.
When tasks can be assigned in any order and multiple tasks may be assigned to
each agent, integer linear programming (ILP) formulations achieve global
optimality but at exponential worst-case cost, rendering them impractical
for large fleets or compute-budgeted settings.

A widely used polynomial-time alternative is the \emph{Sequential Greedy
Algorithm} (SGA)~\cite{minoux1978accelerated}: at each step, the algorithm
selects the (agent, task) pair that offers the greatest marginal gain, assigns
the task, and repeats until no more beneficial assignments remain.
Under utility functions that satisfy the \emph{Diminishing Marginal Gain}
(DMG) property~\cite{fisher1978analysis}, SGA achieves at least half the
optimal assignment value.
This guarantee is tight: there exist instances where SGA achieves exactly
50\% of the optimum.

Centralized methods share two practical weaknesses.
First, they assume that all task and agent information is available at a single
point, which requires either a reliable broadcast channel or a dedicated
coordinator robot.
Both assumptions are fragile in the field: communication failures can prevent
the coordinator from learning about new tasks or updated agent states.
Second, single-point-of-failure: if the coordinator fails, the entire system
halts.
These limitations motivate the decentralized alternatives described next.


%% ── 2.1.2 ──────────────────────────────────────────────────────
\subsection{Decentralized Auction-Based Approaches}
      \label{subsec:bg_decentralized}

Auction-based methods distribute the allocation decision by having agents bid
on tasks and resolve conflicts through a shared consensus protocol.
\citet{gerkey2002sold} and \citet{dias2006market} provide early surveys of
market-based multi-robot coordination, establishing the framework in which
agents bid monetary values on tasks and the highest bidder is awarded the task.
While intuitive, early market-based approaches offered few formal guarantees
on solution quality.

The \emph{Consensus-Based Bundle Algorithm} (CBBA)~\cite{choi2009consensus}
introduced a principled framework that combines greedy bundle building with
a pairwise conflict resolution protocol.
Each agent independently builds a bundle of tasks it would like to execute,
computing bids based on the marginal gain of adding each task to its current
plan.
When agents communicate, they exchange bid and winner vectors and apply a set
of update/reset/leave rules to resolve conflicts.
\citeauthor{choi2009consensus} prove that CBBA converges to a valid
conflict-free assignment (no task assigned to more than one agent) on any
connected graph, and that the assignment achieves at least
50\%-optimality under DMG utilities.

CBBA's key limitation is its \emph{local} consensus protocol: each agent
communicates only with its immediate neighbors, and a single round of exchange
is performed per iteration.
On a graph of diameter $D > 1$, information originating at one node takes $D$
rounds to reach all other nodes.
If agents act on stale information before consensus is reached, conflicting
assignments can persist, and CBBA may converge to an assignment that is not
equivalent to the SGA solution~\cite{kha2025enhancing}.

The \emph{Global Consensus-Based Bundle Algorithm}
(GCBBA)~\cite{kha2025enhancing} corrects this deficiency by performing $2D$
consensus rounds per iteration, where $D$ is the communication graph diameter.
GCBBA also introduces the \emph{Repeated Path Times} (RPT) utility function,
which is provably DMG and better aligned with makespan minimization than the
simple travel-time utility used in CBBA.
Together, these changes allow GCBBA to provably converge to the SGA solution
on any connected graph and with any DMG utility function, while retaining the
decentralized character of CBBA.

GCBBA remains the starting point for this thesis; its handling of disconnected
graphs and dynamic task arrivals is the extension that LCBA provides.


%% ── 2.1.3 ──────────────────────────────────────────────────────
\subsection{Other Decentralized Methods}
      \label{subsec:bg_other_decentralized}

Beyond auction-based methods, several other decentralized approaches have been
proposed for MRTA.
\emph{Distributed Constraint Optimization Problems} (DCOPs) formulate task
allocation as a distributed optimization over coupled agent variables;
algorithms such as Max-Sum and DPOP achieve variable quality-runtime trade-offs
but do not scale well to large numbers of tasks~\cite{chakraa2023optimization}.
\emph{Market-based methods}~\cite{quinton2023mrta} use explicit price dynamics
to drive convergence, offering flexibility at the cost of unpredictable
convergence time.
\emph{Behavior-based methods}~\cite{parker1998alliance} allocate tasks
through emergent competition and reinforcement, with robustness to failures
but without assignment quality guarantees.

The most relevant recent competitor is the \emph{Distributed Matching by Clone
Hungarian-Based Algorithm} (DMCHBA)~\cite{samiei2024dmchba}, which extends
distributed Hungarian matching to multi-robot settings.
DMCHBA achieves strong empirical performance and provides theoretical
guarantees on connected graphs, making it a natural state-of-the-art
baseline for this thesis.
Like CBBA and GCBBA, DMCHBA assumes graph connectivity during each allocation
call.

The common thread across all existing decentralized methods is the assumption
of a \emph{connected} communication graph, at least at the time of allocation.
None of the methods reviewed explicitly handle the case where the graph is
persistently or intermittently disconnected across multiple replanning events.
This gap motivates the LCBA design.


%% ── 2.2 ─────────────────────────────────────────────────────────
\section{Multi-Agent Path Finding}
\label{sec:bg_mapf}

Multi-agent path finding (MAPF) is the problem of computing collision-free
paths for a team of agents from their respective start positions to their goal
positions on a shared graph~\cite{stern2019mapf}.
Formally, a \emph{vertex conflict} occurs when two agents occupy the same node
at the same timestep; an \emph{edge conflict} (or swap conflict) occurs when
two agents traverse the same edge in opposite directions within the same
timestep.
A solution is \emph{valid} if it contains no conflicts.

MAPF is NP-hard to solve optimally~\cite{stern2019mapf}, even for the simplest
cost function (sum of individual path lengths).
Practical approaches therefore trade off optimality against runtime, and the
appropriate trade-off depends heavily on the application: an offline warehouse
layout tool can afford hours of computation for an optimal solution, while an
online robot coordination system must plan within milliseconds per timestep.


%% ── 2.2.1 ──────────────────────────────────────────────────────
\subsection{Classical MAPF Solvers}
\label{subsec:bg_mapf_classical}

\emph{Conflict-Based Search} (CBS)~\cite{sharon2015conflict} is the
most influential optimal MAPF solver.
CBS operates at two levels: a high-level constraint tree (CT) that branches
on conflicts, and a low-level A* that plans each agent's path subject to
accumulated constraints.
CBS is complete and optimal with respect to sum-of-costs and has excellent
average-case performance, though its worst-case complexity is exponential in
the number of conflicts.
CBS remains the algorithm of choice when optimal solutions are needed and
agent counts are small ($n \lesssim 20$ in warehouse grids).

\emph{Windowed Hierarchical Cooperative A*}
(WHCA*)~\cite{silver2005cooperative} is an early suboptimal MAPF method that
applies cooperative A* within a fixed-size time window.
By planning only $w$ timesteps ahead, WHCA* avoids the long-horizon complexity
of full MAPF at the cost of potentially suboptimal paths near the window
boundary.
This idea is formalized and extended by the rolling horizon approach described
in Section~\ref{subsec:bg_mapf_lifelong}.

All A*-based planners share the same admissible heuristic: the shortest
single-agent path from the current cell to the goal, ignoring other
agents~\cite{hart1968astar}.
This heuristic is consistent, ensuring that A* explores nodes in non-decreasing
order of $f$-value and does not revisit nodes.


%% ── 2.2.2 ──────────────────────────────────────────────────────
\subsection{Lifelong and Warehouse MAPF}
\label{subsec:bg_mapf_lifelong}

\emph{Lifelong MAPF}~\cite{li2021lifelong} extends the classical formulation to
settings where agents are never idle: upon completing one goal, each agent is
immediately assigned a new one.
This is the relevant formulation for warehouse automation, where robots
continuously cycle between induct and eject stations.
In the lifelong setting, the natural objective is \emph{throughput} i.e.,
the average number of tasks completed per timestep at steady state(as opposed to 
makespan for a fixed task batch.)

\citet{li2021lifelong} introduce \emph{Rolling Horizon Collision Resolution}
(RHCR) for the lifelong MAPF setting in large warehouses.
RHCR restricts each planning call to a time window of size $w$ and uses
cooperative A* as the inner solver, achieving near-optimal throughput with
computational cost that scales with the window size rather than the total
mission horizon.
RHCR is one of the three planning strategies evaluated in this thesis
(Section~\ref{subsec:method_rhcr}).

A key design insight from \citet{li2021lifelong} is that optimizing for
throughput requires different priority rules than optimizing for makespan.
In particular, agents that are stuck or far from their goals should receive
\emph{higher} planning priority to prevent them from blocking high-throughput
agents near their goals.
This insight informs the priority ordering used in the CA* and PBS planners
in this thesis.


%% ── 2.2.3 ──────────────────────────────────────────────────────
      \subsection{Prioritized Planning and Reservation Tables}
      \label{subsec:bg_prioritized}

      Prioritized planning~\cite{silver2005cooperative} decomposes the joint MAPF
problem into a sequence of single-agent problems.
Agents are assigned a priority order; the highest-priority agent plans its
optimal path first, and each subsequent agent plans a collision-free path
around all higher-priority agents' reservations.

The \emph{reservation table} (also called a space-time reservation table)
stores the space-time positions committed by already-planned agents.
When planning agent $i$'s path, the planner searches the space-time graph
with the reservation table treated as a set of obstacles.

      Prioritized planning is not complete in general: a fixed priority ordering can
create situations where the lowest-priority agent cannot reach its goal because
all routes are permanently blocked by higher-priority agents.
However, \citet{silver2005cooperative} show that on grid graphs without
symmetry constraints, random priority orderings produce very few deadlocks in
practice.
The stuck-detection mechanism in LCBA (Section~\ref{subsec:method_stuck})
provides a practical recovery path for the rare cases where deadlocks do occur.

\emph{Priority Inheritance with Backtracking} (PIBT)~\cite{okumura2019pibt}
      is a reactive variant of prioritized planning that makes greedy decisions
one timestep at a time rather than planning full paths.
Its scalability to large agent counts ($n > 100$) makes it an attractive option
when computational resources are limited, at the cost of providing no
long-horizon path quality guarantee.


%% ── 2.3 ─────────────────────────────────────────────────────────
\section{Integrated Task Allocation and Path Planning}
\label{sec:bg_integrated}

Most MRTA and MAPF research treats the two problems separately.
MRTA methods typically assume that travel times are deterministic and that
paths can always be executed without collision; MAPF methods typically assume
that goals are given and fixed.
In a real multi-robot system, these assumptions are mutually incompatible:
the travel time between two points depends on the paths of other agents (which
in turn depend on their task assignments), and task assignments should account
for the time costs imposed by conflict avoidance.

Tightly integrated approaches, specifically those that solve allocation and path planning
jointly, exist but are computationally demanding.
A small number of works address this integration in the decentralized setting.
\citet{otte2020auctions} study auction-based coordination in communication-limited
environments, but focus on assignment convergence rather than the path planning
layer.
\citet{buckman2019partial} propose partial replanning strategies for dynamic
task allocation, showing that selective replanning of a subset of agents'
paths reduces overhead while maintaining allocation quality.
\citet{chen2021cbbalr} combine CBBA with local path replanning for UAV swarms,
demonstrating that reactive path updates can compensate for stale global plans.

All of these approaches share a key simplifying assumption: the communication
graph is either fully connected at all times, or at most mildly sparse.
None explicitly model the scenario where the graph is \emph{persistently
disconnected} across multiple allocation cycles, as arises when the physical
communication range is shorter than the warehouse diameter.
This thesis addresses exactly this case, combining LCBA allocation
with three interchangeable path planners in a single integrated pipeline.


%% ── 2.4 ─────────────────────────────────────────────────────────
\section{Communication-Constrained Coordination}
\label{sec:bg_communication}

Communication constraints are a recognized challenge in multi-robot systems.
\citet{parker1998alliance} showed that fault-tolerant behavior can emerge from
reactive architectures even without reliable inter-robot communication.
More recently, the MRTA community has studied coordination when communication
is range-limited or intermittent.

\citet{otte2020auctions} analyze auction-based allocation under bandwidth and
range constraints, demonstrating that local auctions can converge to near-global
optima when the communication graph changes slowly.
CBBA was specifically designed for range-limited radios~\cite{choi2009consensus}:
its pairwise consensus protocol requires only one-hop communication and can
operate on any connected graph.
However, CBBA's convergence guarantee breaks down when the graph becomes
disconnected during a consensus round: agents in isolated components may
converge to conflicting assignments that cannot be resolved upon reconnection.

GCBBA~\cite{kha2025enhancing} partially addresses this by noting that
consensus within a connected component is still valid.
However, the original GCBBA formulation does not specify a protocol for
handling reconnection events or for task arrivals that occur during a period
of disconnection.
These gaps are precisely what LCBA fills: by treating each connected component
as an independent allocation instance and triggering reallocation upon topology
changes, LCBA maintains correctness and forward progress regardless of
communication conditions.

The proximity-based communication model used in this thesis is relatively straightforward:
an edge between
agents $i$ and $k$ if and only if their Euclidean distance is at most $\comm$.
 This is a standard abstraction for range-limited radios and is well-suited to
the warehouse grid environment where agents can estimate distances from
grid coordinates.
This model produces graphs that change every timestep as agents move, creating
a highly dynamic topology that no existing MRTA method is explicitly designed
to handle.


%% ── 2.5 ─────────────────────────────────────────────────────────
\section{Warehouse Automation Systems}
\label{sec:bg_warehouse}

Automated warehouse systems have grown dramatically in scale and complexity
over the past two decades.
The Kiva (now Amazon Robotics) system~\cite{wurman2008coordinating} operates
hundreds of drive units on a shared floor, moving shelves to human pickers.
The coordination challenge in such systems is substantial: at peak throughput,
robots must navigate dense environments without collision while continuously
receiving new delivery targets.

Modern warehouse systems are characterized by the \emph{induct--eject}
architecture studied in this thesis: items arrive at a set of induct stations
(conveyor endpoints, receiving docks), are assigned to agents, transported
across the warehouse, and deposited at eject stations (shipping chutes, packing
stations).
Throughput (items processed per unit time) is the dominant performance
metric; makespan is relevant only for finite-batch planning horizons.

The coordination challenges specific to this architecture are well-documented
in the lifelong MAPF literature~\cite{li2021lifelong}.
Deadlocks arise when agents block each other at narrow corridor junctions;
live-lock conditions arise when agents repeatedly defer to each other without
making progress.
Task imbalance defined as some agents receiving disproportionately many or few tasks
degrades throughput and energy efficiency.
These are precisely the failure modes that the allocation and path planning
components of LCBA are designed to avoid.

For the purposes of this thesis, the warehouse is modeled as a discrete grid
with fixed station positions.
While real warehouses are larger and have more complex layouts, the $30 \times
30$ grid used here captures the essential structural features (aisle corridors,
station clusters) while remaining computationally tractable for extensive
parameter sweeps.
Generalisation to larger grids is discussed as a direction for future work in
Chapter~\ref{chap:conclusion}.


%% ── 2.6 ─────────────────────────────────────────────────────────
\section{Positioning of This Work}
\label{sec:bg_positioning}

The review above identifies two partially overlapping lines of work: methods
that integrate task allocation with path planning but assume connected
communication graphs, and methods that handle communication constraints but
treat path planning as a separate concern.
As per our review, no existing work simultaneously addresses all three features that the warehouse
setting demands:

\begin{enumerate}
  \item \textbf{Dynamic task arrivals.}
        Tasks arrive online during the mission rather than being given as a
        fixed batch at $t = 0$.
        This requires event-driven replanning rather than a single allocation
        call.
  \item \textbf{Persistently disconnected communication.}
        The proximity-based communication model produces graphs that may be
        disconnected for sustained periods as robots spread across the warehouse.
        Existing decentralized methods (CBBA, GCBBA, DMCHBA) break down or
        stall when components are isolated.
  \item \textbf{Multiple, interchangeable path planners.}
        Different deployment scenarios (small fleets vs.\ large fleets,
        low-compute hardware vs.\ edge servers) call for different planning
        strategies.
        No existing integrated system provides a modular, planner-agnostic
        allocation layer.
\end{enumerate}

This thesis addresses all three gaps through the LCBA algorithm and its
integration with three path planning modules.
LCBA inherits the 50\%-optimality guarantee and convergence properties
of GCBBA/SGA on connected components, while extending correctness guarantees to
disconnected and dynamically changing topologies.
The empirical evaluation in Chapter~\ref{chap:experiments} quantifies the
benefit of these extensions relative to each of the four baseline methods.
